\section{Conclusion}
\label{sec:conclusion}

This article asked whether the JML inference pipeline of Articles~1--3 can be deployed inside an Agile development workflow without significant overhead or workflow disruption. The answer is yes by construction, and a working CI/CD reference implementation accompanies the article: a CLI driver ({\tt jml-ci}) plus three platform-specific workflow templates (GitHub Actions, GitLab CI, Jenkins), all open source and dogfooded against the inferrer's own repository.

The design rests on three deliberate choices. \emph{Per-PR diff-only analysis with content-hash caching} keeps the per-commit overhead low and the cache invalidation precise enough to survive Article~3's compositional dependencies. \emph{Fail-open semantics} prevents the inference pipeline from blocking the merge train; a developer who is not yet sold on formal methods sees inferrer output as informative annotation rather than gating obstacle. \emph{In-place PR comment updates} keep the review surface clean across multiple pushes. Each choice has a corresponding component in the CLI and the workflow templates, and each is unit-tested.

The empirical evaluation across five external repositories is not yet run; this article specifies it. The pre-registration of the experimental design before the study runs is a deliberate methodological choice: the metrics, decision criteria, and survey instruments are fixed now, so the eventual report cannot retrospectively narrow the criteria to fit favourable outcomes.

This is the fourth article in a sequence: \emph{do specs help LLMs write tests} (Article~1: yes, large effect), \emph{can specs ride alongside compiled artefacts} (Article~2: yes, 100\% fidelity on three OSS libraries), \emph{can specs compose across method boundaries} (Article~3: scaffold demonstrates strict extension over isolated inference), \emph{can the pipeline live inside Agile CI} (Article~4: design + reference implementation, empirical study planned). Together they trace a path from heuristic AST-based inference to a deployable infrastructure.
